\subsection{Hilbert-class primes compatible with the ordinary recursion}

\begin{statement}{A.95}{Proved}
{The required splitting primes and one tower with dense correction image}
Let $M$ be divisible by $24$, and let
$a=3u^8$, $b=-2v^8$, where $u,v$ are positive integers prime
to six and $\gcd(a,b)=1$.
Fix an odd prime $\ell$ and $k\ge1$ with $N=\ell^k\mid M$.
Set
\[
 K=\Q(a^{1/N}),\quad L=K(\mu_N),\quad
 \mathcal B=\Q(\mu_M,b^{1/M}),\quad
 \mathcal R=\Q(\mu_D),\qquad M\mid D.
 \tag{A.95.1}
\]
For any $\rho\in\operatorname{Gal}(\mathcal R/\Q)$ fixing
$\Q(\mu_M)$, and any finite excluded set of rational primes,
the subgroup $C_{\rm split}$ of A.92 is generated by the
split-prime difference classes at primes $q$ outside that set
which simultaneously satisfy
\[
 q\text{ splits completely in }L\text{ and }\mathcal B,
       \qquad \operatorname{Frob}_q|_{\mathcal R}=\rho.
 \tag{A.95.2}
\]
In particular a finite set of such primes supplies all the
required class generators.

There is an ordinary infinite-support pair $(A,B)$ of the
two-prime type of A.44, with $v_3(A)=v_2(B)=1$ and $v_2(A)=0$,
for which the actual nested $3$-power root tower of A.77
satisfies
\[
 \varprojlim_r W_r\longrightarrow C_{n_j}^0
       \text{ is surjective for every }n_j=v_3((j+4)!).
 \tag{A.95.3}
\]
All these maps belong to this one pair and this one tower.
Consequently
\[
 \operatorname{im}\left(\varprojlim_r W_r
                   \longrightarrow\varprojlim_r C_r^0\right)
       \text{ is dense in }\varprojlim_r C_r^0
 \tag{A.95.4}
\]
for its inverse-limit topology.  The construction supplies
finite-support compatible divisor representatives for each
individual target in (A.95.3).
It does not assert that these representatives converge as
the target level increases, or that the class family in A.86
belongs to the image in (A.95.4).
\end{statement}
\begin{proof}
We first prove the simultaneous field statement, including
the effect of the prescribed cyclotomic Frobenius.
Let $H_K/K$ be the primary Hilbert class field and let $F$
be the compositum of $L$ with all $\Q$-conjugates of $H_K$,
as in A.92.  Thus
\[
 F=L\prod_{\sigma\in G_\Q}\sigma(H_K),\qquad
 G=\operatorname{Gal}(F/\Q),\quad
 H=\operatorname{Gal}(F/K),\quad V=\operatorname{Gal}(F/L).
 \tag{A.95.5}
\]
The product contains only finitely many distinct fields.
The maps in A.92 identify $G/V$ with the full affine
holomorph and $H/V$ with $U=(\Z/N\Z)^\times$.
The group $V$ is an abelian $\ell$-group.

\emph{The intersection with the boundary radical.}
The field $F/\Q$ is unramified at two.
Indeed, $a$ is a two-adic unit and $N$ is odd, so adjoining
its $N$th roots is unramified at two: the root algebra over
$\Z_2$ is finite etale, because both the roots and the
derivative $NT^{N-1}$ are units.
The odd cyclotomic extension is unramified there as well.
The extension $F/L$ is unramified at all finite primes by
its Hilbert class field construction.
After any base change, these unramified extensions remain
unramified.  Also $F\mathcal R/\mathcal R$ has $\ell$-power
degree, since $\mathcal R$ contains $\mu_N$ and
$\operatorname{Gal}(F/\Q(\mu_N))$ is an $\ell$-group.

The extension $\mathcal B\mathcal R/\mathcal R
=\mathcal R(b^{1/M})/\mathcal R$ is cyclic by Kummer theory.
Put $s=v_\ell(M)$.
At any prime $\mathfrak p\mid2$ of $\mathcal R$, the normalized
valuation of $b$ is
\[
 v_{\mathfrak p}(b)=e(\mathcal R_{\mathfrak p}/\Q_2),
                  \qquad \ell\nmid v_{\mathfrak p}(b).
 \tag{A.95.6}
\]
The ramification index of a cyclotomic extension at two is
a power of two: its odd-conductor part is unramified and
its two-power conductor part has two-power ramification.
Thus (A.95.6) follows from $v_2(b)=1$.
The local extension obtained by adjoining $b^{1/\ell^s}$
has ramification index $\ell^s$, as its root has valuation
with denominator exactly $\ell^s$.
It is the full $\ell$-primary subextension of the cyclic
Kummer extension, and every nontrivial subextension of
it is ramified at two.
Every subextension of $F\mathcal R/\mathcal R$ is an
unramified $\ell$-extension there.  Therefore the exact
global intersection is
\[
             F\mathcal R\cap\mathcal B\mathcal R
                                      =\mathcal R.
 \tag{A.95.7}
\]

The maximal abelian field of $\mathcal B$ is
$\Q(\mu_M)$, by the private-valuation and quartic-field
calculation in A.44, applied to $b=-2v^8$.
Consequently
$\mathcal B\cap\mathcal R=\Q(\mu_M)$.
The prescribed $\rho$ therefore extends to an automorphism
of $\mathcal B\mathcal R$ which is the identity on
$\mathcal B$.
This verifies the boundary part of the desired Frobenius
before it is combined with (A.95.7).

\emph{The allowed Hilbert Frobenius elements.}
Put $E=F\cap\mathcal R$ and
$V_0=\operatorname{Gal}(F/LE)$.
Since $E/\Q$ is abelian and the maximal abelian field of
$L$ is $\Q(\mu_N)$, one has
$E\cap L=\Q(\mu_N)$.
Restriction of $V$ onto
$\operatorname{Gal}(E/\Q(\mu_N))$ is surjective.
The automorphism $\rho$ fixes $\mu_N$, so the elements
$\sigma\in V$ satisfying $\sigma|_E=\rho|_E$ form a
nonempty coset
\[
        \sigma_0V_0\subset V,
       \qquad V/V_0\text{ has trivial conjugation action by }G.
 \tag{A.95.8}
\]
Triviality of this action follows from the abelian field
$E/\Q$, through the displayed restriction map.
Every $\sigma$ in this coset combines with $\rho$ to give
an automorphism of $F\mathcal R$.
By (A.95.7), it then combines with the boundary automorphism
just constructed to give an automorphism of
$F\mathcal B\mathcal R$.
Chebotarev realizes its conjugacy class at infinitely many
rational primes outside any finite set
\cite[Theorem 8.31]{AgentMilneANT}.
They satisfy (A.95.2), and their Frobenius in $F$ may be
chosen to be $\sigma$ by choosing a prime above them.
All these field and automorphism maps retain the prescribed
cyclotomic restriction.

It remains to show that restricting the Frobenius to this
coset does not remove any generator of $C_{\rm split}$.
Let $C'$ be the subgroup generated by the $R_q$ for the
primes in (A.95.2), and let
$\psi:C\to\Q_\ell/\Z_\ell$ annihilate $C'$.
Pull it back to $H$ by the Artin restriction of A.92 and
write $\lambda=\psi|_V$.
For $g\in G$ put
\[
 \delta_g(v)=\lambda(g^{-1}vg)-\lambda(v),\qquad v\in V.
 \tag{A.95.9}
\]
The split-prime branch formula (A.92.7), for every
$\sigma_0v$ with $v\in V_0$, gives
$\delta_g(\sigma_0v)=0$ for all $g$.
Subtracting its value at $v=1$ proves
$\delta_g|_{V_0}=0$.
Thus every $\delta_g$ factors through $V/V_0$.
This quotient has trivial $G$-action by (A.95.8), so the
usual conjugation identity, evaluated on $v\in V$, gives
\[
  \delta_{g_1g_2}=\delta_{g_1}+\delta_{g_2}
       \quad\text{in }\operatorname{Hom}
              (V/V_0,\Q_\ell/\Z_\ell).
 \tag{A.95.10}
\]
The homomorphism (A.95.10) kills $H$: a character of $H$
is invariant under conjugation by $H$.
It also kills $V$, since $V$ is abelian.
It therefore factors through $(G/V)^{\rm ab}=U$ and is
zero, because $H/V$ maps onto that entire abelianization.
It follows that $\lambda$ is $G$-invariant.
The exact annihilator calculation of A.92 says that
$\psi$ annihilates all of $C_{\rm split}$.
Finite character separation gives $C'=C_{\rm split}$.
This proves the first assertion, including any finite
excluded set and the finite choice of generators.

\emph{Allocating the primes inside one ordinary recursion.}
We now use $\ell=3$ and the literal induction of A.44.
At stage $j$ let $M=M_j=(j+4)!$, let $(a,b)$ be its
current coefficients, and put $n=v_3(M)$.
Before performing the usual step to $M_{j+1}$, form
$K_n=\Q(a^{1/3^n})$ and its Hilbert field $F$ as above.
Choose a cyclotomic modulus $D$ divisible by $M$ and by
sufficiently large powers of every currently controlled
prime, so that multiplication by an integer congruent
to one modulo $D$ preserves all current sum-error bounds.
Take $\rho=1$ in the first assertion.
Choose a finite list of distinct permitted primes
$q_1,\ldots,q_t$ whose difference classes generate
$C_{n,\rm split}$, all outside the coefficient support
and larger than the next finite-field bound.
Replace the first coefficient by
\[
                   a^*=a Q^M,\qquad Q=\prod_{i=1}^t q_i,
                    \qquad b^*=b.
 \tag{A.95.11}
\]
The empty list, and hence $Q=1$, is allowed.

This preliminary operation preserves every invariant
needed for the A.44 step.
It preserves the coefficient shapes $3u^8,-2v^8$, their
coprime supports, and their valuations at two and three.
At a prime other than the $q_i$, an old unit-coordinate
point of $aX^M+bY^M=1$ becomes a unit-coordinate point of
$a^*X^M+bY^M=1$ by replacing $X$ by $X/Q$.
At $q_i$, the element $b$ is an $M$th power and
$a^*$ has positive valuation.  Take $X=1$ and lift an
$M$th root of $(1-a^*)/b$ by Hensel's lemma; its derivative
is a unit since $q_i>M_{j+1}$.
This supplies unit-coordinate points at all these new primes.
The congruence $Q\equiv1\pmod D$ preserves the old
sum-error bounds.  The new primes lie above the bound
at which those bounds were required.
The next A.44 step treats them as old coefficient primes,
using the unit points just constructed.
It chooses its subsequent two new primes outside their
support.  Thus it produces the ordinary next-stage
coefficients with every original invariant intact.

All root fields already fixed at level $3^n$ are unchanged
as embedded number fields: (A.95.11) replaces a chosen
root $\alpha_n$ by the explicit root
$Q^{M/3^n}\alpha_n$.
Every subsequent coefficient change is again a rational
$3^n$th power.
The norm transitions between the fields at successive
root levels are exactly those in (A.86.1).
This verifies that the fields used at different allocation
stages are levels of one fixed final radical tower.

Each allocated prime is never reused as a coefficient
prime and has final ordinary valuation
\[
               v_{q_i}(A)=M,\qquad v_3(v_{q_i}(A))=n.
 \tag{A.95.12}
\]
At level $n$ it splits completely in the fixed field
$L_n$.  Therefore its affine Frobenius acts trivially
on the stable inertia quotient $\Z/3^n\Z$ of A.89.
All $3^n$ residue degrees in that quotient are one,
and the actual projection is an isomorphism
\[
           W_{\infty,q_i}\xrightarrow{\ \sim\ }W_{n,q_i}.
 \tag{A.95.13}
\]
Notice that $q_i$ is still unramified at level $n$;
its total ramification begins above that level.
Equation (A.95.13) uses $n=v_3(M)$ and the actual
inertia quotient, not ramification at the current level.

By A.92, $C_n^0$ is generated by these split-prime
difference images and the images at the primes already
ramified in $L_n$.
The latter also have stable local inverse kernels by
A.89; the prime three contributes zero.
Consequently every element of $C_n^0$ is the image of
a finite sum of local inverse-kernel elements, each
supported at one of these finitely many actual
$A$-support primes.
That finite sum belongs to $\varprojlim W_r$ by the
exact convergence criterion (A.89.4).
This proves (A.95.3) at this stage, and all its supporting
primes and maps persist at every later stage.

Perform this preliminary allocation at every stage and
then the original two-prime step.
The proof of the ordinary limits in A.44 applies verbatim
to the extra finite products: every factor introduced at
stage $j$ has exponent $M_j$, the coefficient supports
remain disjoint, and these exponents kill each fixed
finite multiplicative quotient for all sufficiently
large $j$.
The same sum-error bounds tend to infinity.
Hence there are actual ordinary finite-place limits
$A_p+B_p=1$, with both coordinates nonzero, and the
full completed rational classes have those local values.
Their infinitely many original two-prime factors retain
infinite support.  The real signs remain positive for
$A$ and negative for $B$, as in A.44.
This constructs one pair satisfying all the assertions
above, with $n_j=v_3(M_j)$ tending to infinity.

Finally, take any compatible class family
$c=(c_r)_r\in\varprojlim C_r^0$ and any finite list of
levels up to $r_0$.
Choose $j$ with $n_j\ge r_0$ and use (A.95.3) to choose
$w\in\varprojlim W_r$ with class $c_{n_j}$.
Norm compatibility forces its class to be $c_r$ at
every $r\le r_0$.
Such finite coordinate conditions form a basis for the
inverse-limit topology, proving (A.95.4).
The proof does not control how the finite supports of
these choices change with $j$.  It therefore makes no
claim that the choices converge in the completed divisor
groups, and does not replace density by image membership
in (A.86.5).
\end{proof}

\noindent The finite certificate
\path{code/norm_hilbert_allocation_certificate.py}, with retained
output \path{code/norm_hilbert_allocation_certificate_output.txt},
checks the coset-character argument in explicit permutation-module
extensions, the coprime valuation indices used at two, the stable
root-orbit transitions, and the literal embedded-field power maps.
The simultaneous Chebotarev conditions and the ordinary recursive
construction are proved above.
